\subsection*{Materiales}
\begin{itemize}
	\item 2 Generadores de señales RIGOL, 2 canales
	\item 1 Osciloscópio digital RIGOL
	\item Cables de conexión
	\item Uniones en T
\end{itemize}

\subsection*{Procedimiento}
En la práctica experimental se midió el voltaje de múltiples señales eléctricas,
superpuestas para reproducir las señales anteriormente presentadas, utilizando
el osciloscópio.

Para empezar, se conectaron los dos canales de uno de los generadores a una unión en T
y su salida se conectó al osciloscópio. Se estableció el voltaje pico a pico (VPP) de ambos
canales en $V=\qty{3.000}{\volt}$, y sus frecuencias en $\omega_1=\qty{100.00000}{\hertz}$
y  $\omega_2=\qty{105.00000}{\hertz}$ respectivamente. Se observó la superposición de ambas
señales en el osciloscópio y se registró la curva generada. Luego, manteniendo
el mismo VPP y frecuencia del primer canal se cambió la frecuencia del segundo canal
a $\omega_2=\qty{1300.0000}{\hertz}$. Similarmente, se observó la superposición y se
registró la curva.

Posteriormente, utilizando dos uniones en T adicionales se conectaron dos canales más
provenientes del segundo generador. Manteniendo el VPP como $V=\qty{3.000}{\volt}$, y
la frecuencia $\omega=\qty{100.00000}{\hertz}$ en la Ec. \ref{eq:C}, se generaron los
primeros cuatro términos de la onda cuadrada. Se observó y registró cada una de
las curvas correspondientes a las cuatro sumas parciales. Después, manteniendo la misma
frecuencia y cambiando el VPP a $V=\qty{2.000}{\volt}$ en la Ec. \ref{eq:S}, se generaron los
primeros cuatro términos de la onda diente de sierra. De la misma forma, se observó y registró cada una de las curvas.

Finalmente, las curvas experimentales obtenidas se ajustaron a las Ec. \ref{eq:super},
\ref{eq:C} y \ref{eq:S} respectivamente mediante el método de mínimos cuadrados.

\vspace{1cm}